The treatments of a rating that leaves a cohort come from four literatures: the estimation of rating migrations, survival analysis with censoring and competing risks, the analysis of data that are missing not at random, and the econometrics of partial identification. Table~\ref{tab:related} sets the approaches side by side; the subsections characterise each group and the last states where the proposed method stands.

\begin{table}[htbp]
\centering
\small
\caption{Published treatments of a rating that leaves the cohort during the year, by what each assumes about the ratings that left, what it returns and where it fails when withdrawal is more likely before default; the last row is the proposed method.}
\label{tab:related}
\begin{tabularx}{\textwidth}{p{3.1cm}XXX}
\toprule
Treatment & Assumption on the ratings that left & Output & Failure here \\
\midrule
Removal \citep{schuermann2004} & Independent of default & One PD, $D/(D+S)$ & Understates the PD, and no test of the matrix sees it \\
Keeping in the denominator & None of them would default & One PD, $D/N$ & Understates more than removal \\
Following \citep{eba2017} & Year-end state observed & One PD, $(D+X)/N$ & Needs the state of every rating that left, which a matrix lacks \\
Duration estimator \citep{lando2002} & Independent right censoring & Generator and PD by grade & Biased as removal is, in continuous time \\
Competing-risk bounds \citep{peterson1976} & Dependent risks not identified \citep{tsiatis1975} & Bounds on the net survival function & Bounds wide, no way to place the truth in them \\
Worst case \citep{manski1989} & None & Interval $[D/N,(D+W)/N]$ & Uninformative at typical withdrawal rates \\
Tipping point \citep{cro2020} & A shift, increased until the conclusion changes & Sensitivity curve & No second source to say which shift \\
Proposed anchored method & Odds ratio $k$, read from a second source & PD$(k)$, tipping $k^{*}$, anchored interval & Conditional on the second source defining default alike \\
\bottomrule
\end{tabularx}
\end{table}

\subsection{Withdrawals in the estimation of rating migrations}

The cohort estimator counts, for each state at the start of a period, the share that is in each state at its end \citep{anderson1957}, and it is the estimator behind the transition pages of the agencies. Generator-based estimators place a continuous-time chain behind the annual matrix \citep{jarrow1997}, and the question of whether an empirical matrix has a generator at all has a mathematical answer \citep{israel2001}. Measures of mobility and comparisons of matrices across agencies and periods \citep{jafry2004}, and mixtures of chains that let obligors migrate at different speeds \citep{frydman2008}, show that the one-year matrix is neither the same in every year nor the same for every obligor, which matters for the stability of anything estimated from a pool of years.

Two estimators treat a withdrawal differently. The cohort estimator of \citet{schuermann2004} removes the \wdn{} ratings from the denominator, a treatment they attribute to industry practice, and compares it with a duration-based estimator that uses the time each rating spends in a state; the transition page of the repository of the European Securities and Markets Authority (ESMA) counts the same way. The continuous-time estimator of \citet{lando2002}, implemented in the engine used here, ends the observation of a rating at its withdrawal, which is an assumption of independent right censoring. Confidence intervals for the PDs that these estimators produce are binomial or resampling intervals \citep{schuermann2004,hanson2006}; none of them carries uncertainty about the \wdn{} ratings. A paper on the adjustment of corporate default rates for rating withdrawals \citep{cantor2007} treats the problem as one of its own; the present method differs by stating the PD as a family indexed by the informativeness of withdrawal rather than adopting one adjustment. Guidance for banks asks for the withdrawals to be followed \citep{eba2017}, and work on calibration of PD curves \citep{tasche2013} and on the projection of default rates from one-year matrices \citep{engelmann2024} takes the one-year PD by grade as given.

\subsection{Censoring and competing risks}

Survival analysis has the same problem under another name. The Kaplan-Meier estimator reads the loss of an observation as independent censoring \citep{kaplan1958}. When a rating can end in default or in withdrawal, the two are competing risks, and the net distribution of the time to default, the one that would apply if withdrawal could be switched off, is not identified from the observed data: any joint distribution of the two times with dependent risks has an observationally equivalent one with independent risks \citep{tsiatis1975}. The dependence is therefore an assumption, not a finding. The identified quantity is bounded: \citet{peterson1976} derives bounds for the net survival function given the sub-distribution functions of each cause, and the cumulative incidence of default, the share of ratings that default before withdrawal, is modelled by the sub-distribution hazard of \citet{fine1999}.

These results explain why the cohort PD of a grade is not a PD in the sense the validator needs. The share that defaults before being \wdn{} is the cumulative incidence, and the PD of a rating that is followed to the year end is a net quantity. The bounds of \citet{peterson1976} are the survival-time counterpart of the worst-case bound of the proposed method, and they share its weakness: they bracket the answer, wide, and offer no way to place the truth inside them.

\subsection{Missing not at random and sensitivity analysis}

A \wdn{} rating is an outcome that is missing. Under missingness at random, conditional on what is observed, the missing mechanism can be ignored \citep{rubin1976,little2019}; when the probability of missingness depends on the unobserved outcome, as when withdrawal is more likely before default, ignoring it biases the estimate and the mechanism must be modelled \citep{heckman1979}. The pattern-mixture models of \citet{little1993} give the missing outcomes their own distribution, and the semiparametric models of \citet{scharfstein1999} index the departure from missingness at random by a parameter that is not identified from the data and is varied in a sensitivity analysis.

The practice that the proposed method follows is the sensitivity analysis of clinical trials. Controlled multiple imputation imputes the outcomes of participants who left the study with a shift from the observed arm \citep{carpenter2013}, and the shift is increased until the conclusion of the trial changes; this is the tipping-point analysis, and \citet{cro2020} give a practical guide to it for continuous outcomes. The odds ratio $k$ plays the part of the shift for a binary outcome, and the tipping $k^{*}$ is the shift at which a decision flips. What these analyses lack for a bank is the second source: in a trial there is none, and the sensitivity range is argued from clinical judgement. The repository of the agencies supplies one, imperfectly.

\subsection{Partial identification and intervals for proportions}

The worst-case bound for an outcome that is missing not at random is due to \citet{manski1989}: with no assumption on the missing outcomes, the parameter lies in an interval whose width is the share missing. Bounds under weaker restrictions, with randomised experiments and missing outcomes, follow \citep{horowitz2000}, confidence intervals must cover either the parameter or the whole identified set \citep{imbens2004,chernozhukov2007}, and \citet{tamer2010} surveys the field. The identified set of the PD of a grade under informative withdrawal is the Manski interval, and the sampling interval around its ends is the Jeffreys interval, one of the binomial intervals that has coverage near its nominal level for small counts \citep{brown2001}, alongside the Wilson, Clopper-Pearson and Agresti-Coull intervals \citep{wilson1927,clopper1934,agresti1998}. The proposed method uses the Jeffreys interval, as the engine does for every PD it reports. An interval for a partially identified parameter is built from the intervals of its two ends, which is what the identified band and the anchored interval do; the refinements for the case in which the width of the identified set must itself be estimated are not needed, because here the width is a count that the transition page reports exactly.

\subsection{Where the proposed method stands}

Among the treatments of Table~\ref{tab:related}, three are single answers (removal, keeping, following), one is a continuous-time version of removal, and three are families or bounds that stop short of an answer. The proposed method adds to the last group the one ingredient that turns a sensitivity analysis into an estimate: a second source that identifies $k$. It keeps the sensitivity analysis when no such source exists, because the family PD$(k)$ and the tipping $k^{*}$ need only the transition page, and it reports the anchored interval where one exists. What it takes from the second source is an assumption, the equality of the two sources' definitions of default, and the evaluation of Sections~\ref{sec:results} and~\ref{sec:ablation} measures what the assumption is worth. A bank that records the fate of its own \wdn{} ratings, or an agency with its default file, has a second source of its own and can use the family as the sensitivity analysis and the anchor as the estimate. The application of the tipping-point family to rating withdrawals was not found in the literature searched for this work.
